% content of the chapter relativity
\chapter{Special Relativity}
\label{ch:sr}

Prior to 1905, physics was dominated by Newtonian mechanics and Maxwell’s 
equations of electromagnetism. However, a conflict existed between 
the two. Newtonian physics allowed for infinite speeds, while Maxwell’s 
equations suggested that light travels at a fixed speed, $c$.
The "luminiferous aether" --- a hypothetical medium thought to carry 
light wave --- was proposed to resolve this, but the famous Michelson-Morley
experiment failed to detect it. Albert Einstein resolved this paradox 
not by fixing the old theories, but by reinventing the definitions of space and time.
Principle of absolute time of Section~\ref{rm:book} is rejected.

Special Relativity  is built upon two deceptively simple postulates:
\begin{description}
    \item [The Principle of Relativity] 
       The laws of physics are the same in all inertial frames of reference.
   \item[The Constancy of the Speed of Light] 
       The speed of light in a vacuum ($c \approx 3.00 \times 10^8$ m/s) is constant for all observers, regardless of the motion of the light source or the observer.
\end{description}

If the speed of light must remain constant, then space and time must be flexible to accommodate it.
\begin{description}
    \item[Time Dilation] Time is not absolute. A moving clock ticks more slowly than a stationary clock. As an object approaches the speed of light, time slows down significantly relative to a stationary observer.
 \item[Length Contraction] Objects moving at relativistic speeds appear shorter along the direction of motion to a stationary observer.
 \item[Mass-Energy Equivalence] 
    $E = mc^2$
    This implies that a small amount of mass can be converted into a tremendous amount of energy, the principle behind nuclear power and stars.
\end{description}
General relativity adds time dilation due to gravity to this list.

There is solid  experimental evidence in favor of relativity, both 
from specially designed physical experiments and the facts of everyday life.
The Global Positioning System (GPS) is a daily application of relativity.
 Satellites move fast ($\approx 14,000$ km/h), causing their clocks to tick {\em slower} by 7 microseconds per day according to Special Relativity.
 Satellites are far from Earth's mass (weaker gravity), causing their clocks to tick {\em faster} by 45 microseconds per day according to General Relativity.
Without correcting for this net difference of 38 microseconds daily, GPS accuracy would drift by kilometers every day, rendering the system useless.

The conventional wisdom that classical mechanics is recovered from relativity 
when $v/c \ll 1$ is deeply entrenched in the physics literature. Textbooks 
frequently state that Newtonian formulas emerge by taking the first terms in 
a Taylor expansion of relativistic expressions, or by taking the limit $c \to \infty$.
While such approximations correctly reproduce certain numerical results 
(for one particle), they do not recover the {\em classical mechanics}.
The absence of absolute simultaneity, the lack of global time, incompatibility 
of instantaneous constraints, and the necessity of fields for 
interactions signal a profound structural divergence.

Newtonian mechanics requires a global time parameter shared by all 
degrees of freedom. This time orders events, defines simultaneity, 
and underlies the Hamiltonian formulation. Phase space trajectories 
evolve with respect to a universal time coordinate. The Newtonian $N$-particle 
configuration space $\mathbb{R}^{3N}$ is defined on simultaneous positions of all particles.

Rigid bodies in classical mechanics require instantaneous constraint 
enforcement and a notion of shape preserved by all inertial observers. 
These presuppose global simultaneity. In relativity, Born rigidity 
is possible only under highly restricted motions and cannot be maintained 
for general accelerations. Constraint forces propagate at finite speed, 
so multi-point instantaneous constraints are prohibited.

Constraints defining rigid links, rods, joints, or holonomic relations 
are not relativistically admissible, except as approximations mediated 
through elastic or electromagnetic fields with finite propagation speed. 
The configuration spaces associated with classical mechanical systems 
(e.g., rigid rotors, articulated bodies) have no relativistic analogue.

The famous no-interaction theorem (\cite{CurrieJordanSudarshan1963},\cite{Leutwyler1965}) 
demonstrates that a system of finitely many relativistic particles 
cannot interact through potentials if one demands Poincar'e invariance 
and a single time parameter. Nontrivial interactions require fields. 
Thus classical particle mechanics with inter particle forces has no relativistic extension.

A naive expectation, inherited from Newtonian mechanics, is that a system of
relativistic point particles should admit a Hamiltonian description in which
each particle's canonical position traces out its physical world line, and in
which interactions enter as ordinary potentials depending on the particles'
instantaneous configurations. The \emph{no-interaction theorem} of Currie,
Jordan, and Sudarshan, sharpened by Leutwyler, shows that this expectation
cannot be met: any canonical realization of the Poincar\'e group on a
finite-dimensional many-particle phase space that also respects the natural
requirement that each particle's position be a genuine point on a covariant
world line forces every particle to move on a straight line at constant
velocity, exactly as if it were free. This paper gives a self-contained account
of the theorem---its hypotheses, the mechanism behind its proof, and its
physical origin in the relativity of simultaneity---and then surveys, in a
unified language, the principal strategies that have been developed to obtain
genuinely interacting relativistic dynamics by relaxing exactly one of its
hypotheses: quasi-relativistic (Bakamjian--Thomas) instant-form dynamics,
passage to local field theory, constraint (predictive) mechanics with multiple
mass-shell conditions, manifestly covariant action-at-a-distance without
canonical structure, and the reparametrization-invariant, off-shell dynamics of
Stueckelberg, Horwitz, and Piron. The last of these is treated in some detail
as a case study, since it circumvents the theorem in a particularly transparent
way: by promoting all four components of each particle's spacetime position to
dynamical canonical variables evolving in a fifth, invariant parameter, it
removes the very structure---an external, particle-independent time coordinate
singled out among the dynamical variables---on which the theorem's proof
depends.

This result explains the structural collapse of the Newtonian 
$N$-body problem in relativity. The notion of an instantaneous 
force depending on simultaneous positions is meaningless. 
Relativity allows only local interactions mediated by fields or 
direct interactions (collisions).

One might hope that classical continuum mechanics survives relativity 
even if rigid bodies do not. However, relativistic continuum mechanics 
is fundamentally different from its Newtonian counterpart. 
A relativistic continuum is defined by a stress--energy tensor $T^{ij}$ 
satisfying local conservation laws. There is no concept of a 
perfectly rigid body, and the equations of state must respect causality.

In Newtonian continuum mechanics, pressure and stress can adjust 
instantaneously across the medium. In relativity, signals propagate 
at finite speed, invalidating classical constitutive relations. 
The Newtonian limit of relativistic continua yields compressible, 
not rigid, media. The structural assumptions are thus incompatible.

Given the collapse of constraints, multi-particle structure, 
and continuum rigidity, the only part of classical mechanics 
that survives intact in relativity is the motion of a single 
particle under an external electromagnetic field. 
The Lorentz force law is compatible with Minkowski geometry and 
requires no global time. Maxwell's equations provide the necessary field mediation.
Everything else---interacting particles, rigid bodies, 
constrained systems, Newtonian continua---is not contained within relativity.

Classical mechanics is not a limit case of relativity. 
These are two different models of reality with one particle
problems as common domain and they should cooperate for
physical problems. It is not uncommon that only part of
the system is relativistic, e.g. fast electrons may
interact with slow ions. Besides that for most physical
problems {\em there is} natural frame, the most convenient
for analysis.

\section{Relativistic particle}
\label {sec:sr-relativistic-particle}

 Lagrangian function for relativistic
particle in electromagnetic field can be written as
\cite{LandauLifshitz-2},\cite{GoldsteinPooleSafko2000}:
\begin{equation}
   \label {lagr-3d}
   L=-m c^2 \sqrt{1-\frac{v^2}{c^2}} - e \phi + \frac{e}{c} \bm A \bm v
\end{equation}
So canonical moment $ \bm P = \partial L/\partial \bm v $ is
\begin{equation}
   \label {mon-3d}
    \bm P=\frac{m \bm v}{\sqrt{1-\frac{v^2}{c^2}}} +\frac{e}{c} \bm A 
\end{equation}
and Hamiltonian is
\begin{equation}
   \label {ham-3d}
    H=\sqrt{m^2 c^4+c^2 (\bm P -\frac{e}{c} \bm A)^2} +e \phi
\end{equation}
and equation of motion are
\begin{equation}
   \label {eqm-3d}
   \frac{d \bm P}{d t}=e \bm E + \frac{e}{c} \bm v \times \bm B
\end{equation}

\section{Minkowski spacetime}
\label {sec:sr-minkowski-spacetime}
Newtonian mechanics is inseparable from Galilean spacetime,
which possesses
absolute time and Euclidean space. The geometry consists of a 
three-dimensional spatial manifold evolving in a universal time $t$.
Simultaneity is observer-independent. The symmetry group is 
the Galilei group, comprising translations,rotations and Galilei boost
(see Chapter \ref{ch:cm}):
\begin{equation}
   \label {gal-boost}
   \begin{split}
       t' &= t \\
       \bm x' &= \bm x -\bm v t
   \end{split}
\end{equation}

In contrast, special relativity is formulated in Minkowski spacetime (see \ref{me-cov})
 and the symmetry group is the Poincare group, consisting of
 inhomogeneous linear transformations preserving metric \eqref{minkg}, which also
 contains translations and rotations, but Lorenz boost
 \begin{equation}
    \label {lor-boost}
    \begin{split}
        t' &= \gamma (t-\frac{\bm v \bm x}{c^2}) \\
        \bm x' &= \bm x + \frac{(\gamma -1)}{v^2} (\bm v \cdot \bm x) \bm v-\gamma \bm v t  \\
        \gamma &= \frac{1}{\sqrt{1-\frac{v^2}{c^2}}}
    \end{split}
 \end{equation}
 replaces Galilean one.

Points in that manifold are called events. Taking arbitrary 
event as the origin one can see that
The separation between events (square of invariant interval) 
\begin{equation}
   \label {interval}
    s^2=c^2 t^2-|\bm x|^2
\end{equation}
may be timelike ($s>0$), lightlike ($s=0$) or spacelike ($s<0$).
So when the origin in Minkowski space is fixed, it is divided into
five regions: future ($ s>0 $ $ t>0 $), cone ($ s>0 $ $ t<0 $),
future light cone ($ s=0 $ $ t>0 $), past light cone ($ s=0 $ $ t<0 $)
and elsewhere ($ s<0 $).

The limit $c \to \infty$ is often said to make \eqref{gal-boost} 
from \eqref{lor-boost} 
and collapse the light cones and elsewhere 
into a spatial hyperplane $ t=0 $ --- now. However this limit
does not exists in any mathematical sense. 

There are three types of vectors in Minkowski space,
a vector $v^i$ is timelike if $g(v,v)= g_{ij} v^i v^j=v_i v^i>0$,
lightlike if $v_i v^i=0$ and spacelike if $v_i v^i<0$
Note that for pair of orthogonal vectors only two possibilities exist:
one of them is timelike another is spacelike, or both are lightlike.


It is often said that relativity makes time and space coordinate equal. 
This is plain wrong,
time is not at the least equal to space position, beside difference in sign there is fundamental contrast: we can always stop the particle so
that is coordinates $ \bm x=0 $ but we can not stop time flow. The only documented case \footnote {
Joshua 10:12-14
Then spake Joshua to the LORD in the day when the LORD delivered 
up the Amorites before the children of Israel, and he said in the 
sight of Israel, Sun, stand thou still upon Gibeon; and thou, 
Moon, in the valley of Ajalon.
And the sun stood still, and the moon stayed, until the people 
had avenged themselves upon their enemies. Is not this written 
in the book of Jasher? So the sun stood still in the midst of heaven,
and hasted not to go down about a whole day.} 
raises doubts about its authenticity.

Any  particle of positive mass has its worldline $x^i(s)$ a curve in $M$ 
with timelike tangent vector. Curves with lightlike tangent vector 
correspond to zero mass objects --- light rays, for example.  

Parameter $ s $ may be arbitrary, for timelike curves we choose it to be the
length of the curve.
\begin{equation}
   \label {int4}
   ds^2=c^2 dt^2-(dx^2+dy^2+dz^2)
\end{equation}
$ d \tau = ds/c $ is called proper time.
We define four-velocity $u^i = \frac{dx^i}{ds}$.
\begin{equation}
   \label {vel4}
    u^i=(\gamma,\gamma \bm v/c)
\end{equation}
Note that 
\begin{equation}
   \label {4vnorm}
u_i u^i=1
\end{equation} 
Four-acceleration is defined as 
\begin{equation}
   \label {4acc}
 w^i = \frac{du^i}{ds}  
\end{equation}
From \eqref{4vnorm} we see that $ u_i w^i=0 $.

Four-momentum is defined as $p^i = m c u^i$,hence $ p^i p_i = m^2 c^2$.
Dynamics is encoded by the equation $\frac{dp^i}{d\tau} = f^i$, note that $ u_i f^i = 0$ (mass-shell preservation).

We can formally expand Principle of Least Action on relativity. Luck of global time 
deprives state of the system in fixed time from any physical sense so limits of integration
become undefined. Nevertheless it remains good way of defining equation of motion.

Covariant action my be written as  
\begin{equation}
   \label {act-1}
    S_1=-m c \int (\sqrt{g_{i j} d x^i d x^j}  - \frac{e}{c} A_i d x^i)
\end{equation}
it is evidently independent of parametrization of the worldline. 
Let $ \dot x^i (\lambda)=d x^i / d \lambda $, $ \lambda $ is 
arbitrary. So Lagrangian is:
\begin{equation}
   \label {lag-1}
    L_1=-m c \sqrt{g_{i j} \dot x^i \dot x^j} -\frac{e}{c} A_i \dot x^i
\end{equation}
momentum 
\begin{equation}
   \label {mom-1}
   P_i = -m c \frac{\dot x_i}{\sqrt{\dot x^l \dot x_l}} -\frac{e}{c} A_i
\end{equation}
and equations of motion are
\begin{equation}
   \label {eqm-1}
    m c \frac{\ddot x_j}{\sqrt{\dot x^l \dot x_l}} (\delta_i^j -
    \frac{\dot x^i \dot x_j}{\dot x^l \dot x_l}) = \frac{e}{c} F_{j i} \dot x^i
\end{equation}
The equation \eqref{eqm-1} does not include longitudinal component of 
4-velocity. If we choose it to satisfy \eqref{4vnorm} (i.e. use natural
parametrisation we get:
\begin{equation}
   \label {eqm-2}
    m c w_j = \frac{e}{c} F_{j i} u^i
\end{equation}
which corresponds to momenta:
\begin{equation}
   \label {mom-2}
   P_i = -m c  u_i -\frac{e}{c} A_i
\end{equation}
and Lagrangian:
\begin{equation}
   \label {lag-2}
    L_2=- \frac{m c}{2} g_{i j} \dot x^i \dot x^j -\frac{e}{c} A_i \dot x^i
\end{equation}
We ignore restriction \eqref{4vnorm} calculating variation of action. 
Though it often said that \eqref{4vnorm} is constraint it is not
actually, it is constant of motion and should be added to \eqref{eqm-2}
as restriction on initial conditions, to provide physically meaningful
solutions. The same situation exactly as with Lorenz gauge for wave
equation \eqref{me-potc} and reduction in classical mechanics.

\section{Accelerated Motion and Bell's Paradox}
\label{sec:sr-bell}
Bell's spaceship paradox provides a particularly sharp and conceptually
transparent demonstration of the incompatibility between rigidity and
relativistic kinematics. Consider two identical spaceships, initially at rest
in an inertial frame $S$, separated by a fixed distance $L$. At time $t=0$ (as
measured in $S$), both spaceships begin to accelerate identically, with the
same prescribed acceleration history in the coordinates of $S$, so that their
spatial separation in $S$ remains constant and equal to $L$ at all times.

From the perspective of the inertial frame $S$, nothing unusual appears to
occur: the ships remain parallel worldlines with fixed separation while their
velocities increase synchronously. However, this description already hides a
crucial relativistic assumption—namely, that ``distance at a given time'' is
being defined using the simultaneity convention of $S$, which has no invariant
meaning.

To assess the physical state of the system, consider instead the instantaneous
rest frame $S'$ of the ships at some later time, when their common velocity is
$v$. Events that are simultaneous in $S$ are not simultaneous in $S'$, and the
spatial distance between the ships measured in $S'$ is obtained by Lorentz
transformation. One finds
\begin{equation}
L' = \gamma(v)\,L,
\end{equation}
where $\gamma(v) = (1 - v^2/c^2)^{-1/2}$. Thus the proper distance between the
ships increases monotonically as their velocity increases.

If the two ships are connected by a string or rod, this increasing proper
separation manifests itself as tensile stress. Since no real material can
sustain unbounded stress, the string must eventually break. Importantly, this
failure occurs despite the fact that the coordinate distance in $S$ is
constant. The paradox therefore exposes the inadequacy of Newtonian intuition,
according to which identical accelerations preserve internal distances.

The resolution lies in the distinction between coordinate acceleration and
proper acceleration. Identical coordinate accelerations do not correspond to
identical proper accelerations, nor do they preserve Born rigidity. In fact,
maintaining constant proper distance would require a precisely tuned,
position-dependent acceleration field, in which the trailing spaceship
accelerates more strongly than the leading one. Such acceleration profiles are
highly constrained and cannot be arbitrarily prescribed.

Bell's spaceship paradox thus demonstrates that even the simplest attempt to
impose rigidity under acceleration leads inevitably to internal stress and
material failure. The lesson is not merely dynamical but conceptual: there
exists no observer-independent notion of a rigidly accelerating extended body.
Any attempt to enforce rigidity necessarily privileges a particular foliation
of spacetime and thereby violates the spirit of relativistic covariance.

\section{No-interaction theorem}
\label{sec:sr-no-int}

Non-relativistic $N$-body mechanics tolerates interaction with remarkable ease.
Given canonical variables $({\bm q}_i,{\bm p}_i)$, $i=1,\dots,N$, essentially any
potential $V({\bm q}_1,\dots,{\bm q}_N)$ may be added to the free Hamiltonian
$H_0=\sum_i {\bm p}_i^2/2m_i$, and the resulting dynamics remains invariant under
the full Galilei group: translations, rotations, and boosts all act on the
${\bm q}_i$ in the same simple, interaction-independent way, because Galilean time
is absolute and a ``simultaneous configuration'' $({\bm q}_1(t),\dots,{\bm q}_N(t))$
means the same thing to every inertial observer.

It is tempting to expect the same freedom in special relativity: replace $H_0$
by $\sum_i\sqrt{{\bm p}_i^2+m_i^2}$, add an interaction term, and require the ten
generators of the Poincar\'e group to close under Poisson brackets. This
expectation fails, and it fails for a deep and structural reason rather than a
technical one. The failure is codified in the \emph{no-interaction theorem},
established by Currie, Jordan, and Sudarshan in 1963 \cite{CurrieJordanSudarshan1963}
and given a shorter, more transparent proof by Leutwyler in 1965
\cite{Leutwyler1965}. Informally: \emph{any} canonical,
Poincar\'e-invariant Hamiltonian description of a finite number of particles,
in which each particle's canonical position is required to be literally a point
on that particle's own covariant world line, describes only free motion---no
scattering, no bound states, no force of any kind.

This is a genuinely surprising conclusion, and it is important to be precise
about what does and does not cause it. It is \emph{not} a statement that
special relativity forbids interacting systems; nature is full of them. It is a
statement about the incompatibility of three specific technical requirements
when combined in the most naive way one could imagine writing down an $N$-body
relativistic Hamiltonian theory: (i) canonical Poisson-bracket structure on a
phase space built from finitely many particle variables, (ii) an \emph{exact}
realization of the Poincar\'e Lie algebra on that phase space, and (iii) the
\emph{world-line condition}, which demands that each particle's canonical
position transform under boosts exactly as the spatial part of a single
relativistic world line would. Every known scheme for describing interacting
relativistic particles---quantum field theory, constraint (``predictive'')
mechanics, the practical Hamiltonian dynamics used in relativistic few-body and
nuclear physics, manifestly covariant action-at-a-distance theories, and
reparametrization-invariant mechanics of the Stueckelberg--Horwitz--Piron
type---proceeds by identifying and giving up exactly one of these three legs.
None of them ``beats'' the theorem; they sidestep it by changing the question.


In this section, units are chosen so that $c=1$, the metric signature is
$(+,-,-,-)$, Greek indices $\mu,\nu=0,1,2,3$ run over spacetime and Latin
indices $a,b,k,l=1,2,3$ over space, and repeated spatial indices are summed
whether raised or lowered (their placement then carries no meaning). Particle
labels are $i,j=1,\dots,N$.

\subsection{The Hamiltonian Framework for Relativistic Particle Mechanics}

Dirac observed in 1949 \cite{Dirac1949} that a Poincar\'e-invariant Hamiltonian
theory can be built around any of several choices of ``instant''---a family of
spacelike (or null) hypersurfaces on which initial data are specified---and
that the choice determines which of the ten Poincar\'e generators are
\emph{kinematic} (interaction-free, fixed by the choice of instant alone) and
which are \emph{dynamical} (carrying the interaction). The three forms in
common use are:
\begin{description}
    \item [Instant form.] Data are specified on the hyperplane
        $x^0=t=\text{const}$. Spatial translations and rotations map such
        hyperplanes into themselves, so ${\bm P}$ and ${\bm J}$ are kinematic; boosts
        and time translations mix $t$ with the other coordinates, so $H\equiv
        P^0$ and ${\bm K}$ (the boost generators) are dynamical.
    \item [Point form.] Data are specified on a mass hyperboloid $x^\mu
        x_\mu=\text{const}$. The full homogeneous Lorentz subgroup (${\bm J}$ and
        ${\bm K}$) is then kinematic, while all four components of $P^\mu$ are
        dynamical.
    \item [Front form.] Data are specified on a null hyperplane such as
        $x^0+x^3=\text{const}$. The seven-parameter subgroup stabilizing this
        plane is kinematic; the remaining three generators are dynamical.
\end{description}

The no-interaction theorem was originally proved in the instant form and later
extended to the point and front forms; the obstruction is the same in all
three, since it concerns the world-line condition rather than any peculiarity
of a particular choice of instant. We work exclusively in the instant form
below, both because it is the most intuitive and because it is the form in
which the theorem's content---and the strategies for evading it---are easiest
to state.

In the instant form, each particle $i$ is described at a common time $t$ by a
canonical pair $({\bm q}_i(t),{\bm p}_i(t))\in\mathbb{R}^3\times\mathbb{R}^3$, with
\begin{equation}
\{q_i^a,p_j^b\}=\delta_{ij}\delta^{ab}, \qquad \{q_i^a,q_j^b\}=0=\{p_i^a,p_j^b\}.
\end{equation}
The generators $H,{\bm P},{\bm J},{\bm K}$ are functions on this phase space (with $t$ an
external, non-dynamical parameter) required to satisfy the Poincar\'e algebra
\begin{gather}
\{J^a,J^b\}=\epsilon^{abc}J^c, \quad \{J^a,K^b\}=\epsilon^{abc}K^c, 
    \quad \{K^a,K^b\}=-\epsilon^{abc}J^c, \label{eq:algebra1}\\
\{J^a,P^b\}=\epsilon^{abc}P^c, \quad \{J^a,H\}=0, \quad \{K^a,H\}=P^a,
    \quad \{K^a,P^b\}=\delta^{ab}H, \quad \{P^a,P^b\}=0. \label{eq:algebra2}
\end{gather}

\paragraph{The free particle.} For a single free particle with
$H_0=\omega\equiv\sqrt{{\bm p}^2+m^2}$, the boost generator is
\begin{equation}
K_0^a = \omega\, q^a - t\,p^a .
\label{eq:K0}
\end{equation}
It is a useful and instructive exercise---carried out here explicitly, since it
fixes the sign conventions used throughout the paper and makes the later
definition of the world-line condition transparent rather than asserted---to
check \eqref{eq:K0} against \eqref{eq:algebra2}. Using
$\{q^a,f({\bm p})\}=\partial f/\partial p_a$ and the Leibniz rule
$\{fg,h\}=f\{g,h\}+g\{f,h\}$:
\begin{align}
\{K_0^a,H_0\} &= \{\omega q^a,\omega\} - t\{p^a,\omega\} = \omega\{q^a,omega\}
    = \omega\cdot\frac{p^a}{\omega} = p^a = P^a, \\
\{K_0^a,P^b\} &= \{\omega q^a,p^b\} = \omega\,\delta^{ab} = \delta^{ab}H_0, \\
\{K_0^a,K_0^b\} &= \{\omega q^a,\omega q^b\} = q^bp^a - q^ap^b = -\epsilon^{abc}J^c,
\end{align}
with ${\bm J}={\bm q}\times{\bm p}$ in the last line, as required. (The first equality in
the third line uses $\{\omega q^a,\omega q^b\}=q^bp^a-q^ap^b$, obtained by
the same Leibniz expansion together with $\{q^a,\omega\}=p^a/\omega$; the
$t$-dependent cross terms $-t\{\omega q^a,p^b\}-t\{p^a,\omega q^b\}$ cancel
exactly.) Hamilton's equations then give $\dot q^a=\{q^a,H_0\}=p^a/\omega$,
the correct relativistic velocity, so that $x^\mu(t)=(t,{\bm q}(t))$ is indeed the
free particle's world line.

Now compute how the free boost generator acts on the position itself:
\begin{equation}
\{K_0^a,q^b\} = \{\omega q^a,q^b\} = q^a\{\omega,q^b\} = -q^a\frac{p^b}{\omega},
\end{equation}
so that, restoring the explicit $t$-dependence of $K_0^a$,
\begin{equation}
\{K_0^a,q^b\} = t\,\delta^{ab} - q^a\,\dot q^b .
\label{eq:wlc-free}
\end{equation}
Equation~\eqref{eq:wlc-free} says something geometrically simple: an
infinitesimal boost moves the snapshot ${\bm q}(t)$ exactly the way it would move
if one boosted the actual four-dimensional world line $x^\mu(\tau)$ and re-read
off its spatial position at the (now different) boosted time. The term
$t\delta^{ab}$ is the ``kinematic'' piece present even for a particle at rest;
the term $-q^a\dot q^b$ is what converts the naive boost of a snapshot into the
boost of a genuine trajectory, and it does so using \emph{that particle's own}
instantaneous velocity, and only that particle's velocity.

\begin{definition}[World-line condition]
A choice of interacting generators $H,{\bm K}$ satisfies the world-line condition
    (WLC) if, for every particle $i$,
\begin{equation}
\{K^a,q_i^b\} = t\,\delta^{ab} - q_i^a\,\dot q_i^b, \qquad \dot q_i^b \equiv \{q_i^b,H\}.
\label{eq:wlc}
\end{equation}
\end{definition}

The WLC is the precise sense in which $q_i(t)$ is required to be ``the position
of particle $i$ at time $t$'' in a fully relativistic, frame-independent sense:
it guarantees that the curve $t\mapsto q_i(t)$ is, to every order in the boost
parameter, the trace on the hyperplane $x^0=t$ of one and the same covariant
world line, whatever the interaction turns out to be. It is automatically
satisfied by the free system \eqref{eq:K0}, and it is the natural---indeed,
seemingly unavoidable---requirement one would impose if asked to write down a
relativistic generalization of Newtonian $N$-body mechanics with each
particle's trajectory kept as a directly meaningful physical object.

\begin{theorem}[Currie--Jordan--Sudarshan 1963; Leutwyler 1965]
\label{thm:noint}
Let $N\ge 2$ and let $({\bm q}_i,{\bm p}_i)_{i=1}^N$ be canonical variables as above.
    Suppose $H,\ {\bm P}=\sum_i{\bm p}_i,\ {\bm J}=\sum_i{\bm q}_i\times{\bm p}_i,\ {\bm K}$ satisfy
    the Poincar\'e algebra \eqref{eq:algebra1}--\eqref{eq:algebra2} and the
    world-line condition \eqref{eq:wlc} for every particle $i$. Then
\begin{equation}
\dot{{\bm p}}_i \;=\; \{{\bm p}_i,H\} \;=\; 0 \qquad (i=1,\dots,N),
\end{equation}
i.e.\ every particle moves on a straight timelike line at constant velocity:
    the dynamics is, trajectory for trajectory, indistinguishable from that of
    $N$ free particles.
\end{theorem}

The force of the theorem is that it applies however $H$ is chosen: one is not
restricted to $H=\sum_i\omega_i+V$ for some innocuous-looking $V$ and then told
that $V$ must vanish; \emph{any} $H$ consistent with
\eqref{eq:algebra1}--\eqref{eq:algebra2} and \eqref{eq:wlc} yields force-free
motion for every particle. An interaction may still shift the numerical value
of the total energy in a way that is not simply additive across particles, but
it cannot bend a single trajectory. This is what makes the result a genuine
\emph{theorem} about relativistic mechanics rather than a statement about the
unsuitability of one particular ansatz for $V$.


The complete proof  is a self-contained exercise in repeatedly applying the
Jacobi identity to \eqref{eq:algebra1}--\eqref{eq:algebra2} together with
\eqref{eq:wlc}.  The logical skeleton is as follows.

\begin{enumerate}
\item Write $H=H_0+V$ and ${\bm K}={\bm K}_0+{\bm W}$, where $H_0=\sum_i\omega_i$ and
    ${\bm K}_0=\sum_i(\omega_i{\bm q}_i-t{\bm p}_i)$ are the free generators of
        Section~\ref{sec:wlc}, and $V,{\bm W}$ are the (a priori unknown)
        interaction-dependent corrections. Since ${\bm p}$ and ${\bm J}$ are required
        to remain purely kinematic, all of the interaction is confined to $V$
        and ${\bm W}$.
\item The relation $\{K^a,H\}=P^a$ in \eqref{eq:algebra2}, expanded to first
    order in the interaction and combined with the free identity
        $\{K_0^a,H_0\}=P^a$, yields a linear first-order partial differential
        equation relating ${\bm W}$ to the derivatives of $V$.
\item The world-line condition \eqref{eq:wlc}, applied to the interacting
    ${\bm K}$, fixes the dependence of ${\bm W}$ on $V$ still further: because
        \eqref{eq:wlc} must hold with the \emph{interacting} velocity $\dot
        q_i^b=\{q_i^b,H\}$ on its right-hand side, ${\bm W}$ cannot be chosen
        independently of $V$ once \eqref{eq:wlc} is imposed exactly.
\item The remaining algebra relation $\{K^a,K^b\}=-\epsilon^{abc}J^c$,
    together with ${\bm J}$ being purely kinematic, supplies a second, independent
        PDE constraint on ${\bm W}$.
\item Steps 2--4, closed under the Jacobi identity, form an overdetermined
    system of first-order PDEs for $V$ in the particle positions and momenta.
        Currie's original argument solves this system directly for $N=2$
        ; Leutwyler's shorter argument treats general $N$ at
        once \cite{Leutwyler1965}. In both cases, the only solutions compatible
        with translation invariance, rotation invariance, and the physically
        necessary boundary condition that well-separated particles decouple,
        have $\partial V/\partial q_i^a\equiv 0$ for every $i$ and $a$---i.e.,
        $V$ can depend on the momenta alone, and even then only in a way that
        leaves $\{{\bm p}_i,H\}=-\partial V/\partial{\bm q}_i=0$ identically along
        every trajectory.
\end{enumerate}

The upshot is that the world-line condition, which looks like an innocuous
statement about how positions transform, is in fact so rigid that it leaves no
room for the boost generator to be deformed by an interaction in any way that
survives contact with the rest of the algebra---except in the trivial direction
that does not affect any particle's acceleration.

The mechanism behind the theorem has a transparent physical reading. In
Newtonian mechanics, ``the configuration of the system at time $t$,'' i.e.\ the
set of events $\{(t,{\bm q}_i(t))\}_{i=1}^N$, is a frame-independent notion: every
observer agrees on which events are simultaneous, so an instantaneous force law
$F_i({\bm q}_1,\dots,{\bm q}_N)$ that depends on that common configuration is
automatically consistent for all observers.

In special relativity, simultaneity is frame-dependent. If particle $j$ is to
influence particle $i$'s trajectory at the event $(t,{\bm q}_i(t))$ through some
notion of ``the state of particle $j$ at the same time,'' different inertial
observers will disagree about which point on particle $j$'s world line counts
as simultaneous with that event. A genuinely covariant equal-time interaction
law would therefore have to look different---carry different information into
$\dot{{\bm p}}_i$---in different frames, unless the interaction has no dynamical
effect at all. The no-interaction theorem is the precise Hamiltonian-mechanics
statement of this fact: it shows that no choice of ${\bm K}$ can simultaneously (a)
generate the correct boost transformation of the total system, (b) treat each
$q_i$ as a literal point on its own covariant world line, and (c) let
${\bm p}_j(t)$ feed into $\dot{{\bm p}}_i(t)$ at all. Something has to give, and the
theorem tells us that, within this framework, what gives is the interaction
itself.


Theorem~\ref{thm:noint} rests on three ingredients: (I) a finite-dimensional
canonical phase space built purely from particle variables, (II) an
\emph{exact} realization of the Poincar\'e algebra on that phase space, and
(III) the world-line condition \eqref{eq:wlc}. Every workaround below relaxes
exactly one of these.

\paragraph{Relaxing the world-line condition: quasi-relativistic (Bakamjian--Thomas) dynamics}

The most direct response is to keep (I) and (II) and drop (III): allow ${\bm
q}_i(t)$ to be merely a convenient canonical label for particle $i$, not a
claim that ${\bm q}_i(t)$ is literally a point on a covariant world line.
Bakamjian and Thomas showed in 1953  (building on earlier work of Thomas  and
later extended by Foldy ) how to do this constructively. Start from the free
system and pass, by an ordinary canonical (contact) transformation, to
center-of-mass variables $({\bm R},{\bm p})$ and a set of internal (relative)
variables $(\boldsymbol\xi_a,\boldsymbol\pi_a)$, in terms of which the free
invariant mass is
$M_0(\boldsymbol\xi,\boldsymbol\pi)=\sum_i\omega_i^{\text{int}}$. Interaction
is then introduced by replacing $M_0$ with
\begin{equation}
M = M_0(\boldsymbol\xi,\boldsymbol\pi) + U(\boldsymbol\xi,\boldsymbol\pi), 
   \qquad \{U,{\bm R}\}=\{U,{\bm P}\}=0,
\end{equation}
for an \emph{arbitrary} internal potential $U$ subject only to translation and
rotation invariance, and setting
\begin{equation}
H=\sqrt{{\bm p}^2+M^2}, \qquad {\bm K} = -{\bm R}\,H + (\text{spin-dependent terms}).
\end{equation}
One checks that this reproduces the full Poincar\'e algebra
\eqref{eq:algebra1}--\eqref{eq:algebra2} exactly, for any admissible $U$: the
construction is essentially guaranteed to work because it never touches ${\bm
p},{\bm J}$ and builds $H,{\bm K}$ from a single interacting scalar $M$ in
exactly the functional pattern that made the free case close.

The price is precisely the WLC. Once $U\neq0$, the physical positions ${\bm
q}_i$ recovered from $({\bm R},\boldsymbol\xi)$ no longer satisfy
\eqref{eq:wlc}: boosting the system mixes center-of-mass and internal variables
in a way that does not correspond to boosting each ${\bm q}_i(t)$ as an
independent covariant trajectory. The individual ``positions'' become
interaction-dependent, frame-relative bookkeeping devices, while the physically
observable content---the total four-momentum, the invariant mass spectrum, and
the $S$-matrix---remains exactly Poincar\'e invariant. This is not a defect so
much as the whole point: it is precisely the class of construction underlying
most of the practical relativistic Hamiltonian dynamics used in nuclear and
hadronic few-body physics (see the review of Keister and Polyzou).

A second response keeps (II) and (III) in a suitably generalized sense and
drops (I): abandon the idea of a \emph{finite-dimensional, closed} particle
phase space altogether. Introduce a field $\phi(x)$ (or $A^\mu(x)$, or a
non-Abelian gauge field) as an independent canonical system with its own,
infinite-dimensional phase space, and let the particles couple to it locally,
e.g.\ through a Lagrangian density built by adding the particles' world-line
actions to $\mathcal L_{\text{field}}(\phi,\partial\phi)$ and a local coupling
term. The ten Poincar\'e generators are then given by the Noether charges of
the total (particle + field) stress-energy tensor,
\begin{equation}
P^\mu=\int d^3x\, T^{0\mu}, \qquad J^{\mu\nu}=
    \int d^3x\,\big(x^\mu T^{0\nu}-x^\nu T^{0\mu}\big),
\end{equation}
and close the algebra automatically for any Poincar\'e-invariant, local
Lagrangian density---no analogue of the CJS obstruction arises, because the
hypothesis (I) simply does not describe this theory. The particles are now
\emph{open} subsystems, continuously exchanging four-momentum with the field;
only the combined system is conserved and covariant, and there is no
equal-time, direct-action potential $V({\bm q}_1,\dots,{\bm q}_N)$ for the
theorem to constrain in the first place.

This is, historically and practically, the dominant resolution: it is why
interactions in particle physics are described by field exchange rather than
direct action-at-a-distance potentials. The cost is equally well known:
infinitely many degrees of freedom, the need for renormalization at the quantum
level, and, already at the classical level, self-energy divergences for point
sources unless the theory is regularized (the classical electron self-force
problem is the oldest example).

A third response keeps (I) and (III) in spirit but drops the requirement,
hidden inside (II), that there be a single external time $t$ common to all
particles and a single Hamiltonian $H$ generating evolution in it.
Instead---following Todorov , Komar , and Droz-Vincent ---one imposes, for each
particle, a mass-shell-type constraint
\begin{equation}
\phi_i \;\equiv\; p_i^2 - m_i^2 - \Phi_i(x_1,\dots,x_N,p_1,\dots,p_N) \;\approx\;
    0, \qquad i=1,\dots,N,
\label{eq:constraint}
\end{equation}
treated \`a la Dirac as a first-class constraint generating evolution of
particle $i$ along its own parameter $\tau_i$, subject to the compatibility
(``Poisson commutativity'') conditions $\{\phi_i,\phi_j\}\approx 0$ for all
pairs, which ensure that the various $\tau_i$-flows do not obstruct one another
on the constraint surface. Poincar\'e invariance is imposed directly on the
$\phi_i$ (each must be built from Lorentz scalars, typically the invariant
separation of $x_i-x_j$ projected orthogonally to the total momentum), rather
than through a WLC tied to a single $t$.

Because there is no single global $H$ and no single $t$, hypothesis (II) as
used in Theorem~\ref{thm:noint} is simply not the structure this formalism has,
and the CJS/Leutwyler obstruction has nowhere to act. This approach yields
genuinely interacting, exactly Poincar\'e-covariant two- and many-body dynamics
and underlies a substantial body of work on relativistic bound-state problems,
including the two-body Dirac equations of Crater and Van Alstine used in
relativistic quark models \cite{CraterVanAlstine}. The price is technical: the
compatibility conditions on the $\phi_i$ significantly restrict the admissible
form of the interaction functions $\Phi_i$, and the resulting formalism is
considerably more elaborate than a single Hamiltonian.

A fourth response abandons canonical/Hamiltonian structure altogether and works
directly with a manifestly Poincar\'e-invariant multi-time \emph{action}, built
entirely from scalars. The paradigm example is the Wheeler--Feynman
electrodynamics of direct interparticle action , with Fokker-type action
\begin{equation}
S = -\sum_i m_i\int d\tau_i\,\sqrt{-\dot x_i\cdot\dot x_i}\; 
    -\;\frac12\sum_{i\ne j} e_ie_j \iint d\tau_i\,d\tau_j\;
    \delta\big((x_i-x_j)^2\big)\;\dot x_i\cdot\dot x_j .
\end{equation}
This action is trivially Poincar\'e invariant (it is built from the invariants
$\dot x_i^2$ and $(x_i-x_j)^2$) and genuinely interacting. It evades the
theorem because the theorem's hypotheses are stated for a \emph{canonical
Hamiltonian} formulation with a single common time; a covariant, multi-time
Lagrangian/action formulation is simply outside that hypothesis space.

The catch is that this elegance does not survive an attempt to Hamiltonize the
theory on an equal-time slice: doing so either reproduces the CJS/Leutwyler
obstruction directly, or requires retarded (and advanced) interactions that
make the Cauchy data effectively infinite-dimensional---in effect smuggling a
field theory back in through the back door (Wheeler and Feynman's own
resolution, via the absorber, is precisely of this character). Historically,
understanding why the beautiful Fokker-type action program resisted a clean
particle-only Hamiltonization was one of the motivations for the original CJS
investigation; the later, purely Lagrangian proof of the theorem by Marmo,
Mukunda, and Sudarshan  closes the loop by showing that the obstruction can be
seen without ever passing through a Hamiltonian formulation at all.

A fifth response, treated in detail in Section~\ref{sec:shp}, keeps (I) and
(II) but removes the very feature of (III) that makes it so restrictive: the
identification of the evolution parameter $t$ with one component, $x_i^0$, of a
particle's own position. In the Stueckelberg--Horwitz--Piron (SHP) formalism
\cite{Horwitz2015}, evolution is generated by a single invariant parameter
$\tau$, distinct from every $x_i^\mu$, and \emph{all four} components of each
particle's position become dynamical canonical variables. As shown below, this
removes the inhomogeneous $t\,\delta^{ab}$ structure that drives the CJS
argument, and a fully interacting, exactly Poincar\'e-covariant many-body
Hamiltonian becomes possible.
\subsection{ Stueckelberg--Horwitz--Piron Dynamics}

Stueckelberg introduced, in 1941, a reparametrization-invariant single-particle
mechanics in which the particle's four-momentum is not required to satisfy the
mass-shell relation at every instant , motivated originally by the problem of
describing pair creation and annihilation classically. Horwitz and Piron
generalized this to a consistent many-body theory in 1973 ; a systematic modern
account is given in Horwitz's monograph \cite{Horwitz2015}.

Introduce a single real parameter $\tau$, required to be a genuine Poincar\'e
\emph{scalar}---unaffected by translations, rotations, or boosts---and playing,
for the whole $N$-body system, exactly the role that absolute time plays in
Newtonian mechanics. Each particle $i$ is described by \emph{four} independent
canonical position variables $x_i^\mu(\tau)$ together with conjugate momenta
$p_{i\mu}(\tau)$,
\begin{equation}
\{x_i^\mu,p_{j\nu}\} = \delta_{ij}\,\delta^\mu_\nu, \qquad \{x_i^\mu,x_j^\nu\}=0=\{p_{i\mu},p_{j\nu}\}.
\end{equation}
Note the crucial structural difference from Section~\ref{sec:wlc}:
$x_i^0(\tau)$ is now a dynamical variable on exactly the same footing as
$x_i^1,x_i^2,x_i^3$, not an external parameter equated with $\tau$. Evolution
is generated by a single scalar Hamiltonian $K(x_1,\dots,x_N,p_1,\dots,p_N)$,
\begin{equation}
\frac{dx_i^\mu}{d\tau} = \frac{\partial K}{\partial p_{i\mu}}, \qquad \frac{dp_{i\mu}}{d\tau} = -\frac{\partial K}{\partial x_i^\mu},
\end{equation}
and a natural choice, directly analogous to a nonrelativistic $N$-body Hamiltonian, is
\begin{equation}
K = \sum_i \frac{p_i^\mu p_{i\mu}}{2M_i} + V\big((x_i-x_j)^2,\dots\big),
\label{eq:shpK}
\end{equation}
with $V$ built from Lorentz invariants of the particle separations (and
possibly of the momenta). No mass-shell constraint $p_i^2=m_i^2$ is imposed at
each $\tau$: $p_i^2(\tau)$ is a dynamical quantity, generally not constant
along the interacting evolution---particles are, in this sense, \emph{off
shell}. The ten Poincar\'e generators
\begin{equation}
P^\mu=\sum_i p_i^\mu, \qquad J^{\mu\nu}=\sum_i\big(x_i^\mu p_i^\nu - x_i^\nu p_i^\mu\big)
\end{equation}
close the Poincar\'e algebra \emph{exactly}, for \emph{any} choice of invariant
potential $V$ in \eqref{eq:shpK}---no restriction on $V$ analogous to the
conclusion of Theorem~\ref{thm:noint} arises.


It is worth being explicit about exactly which structural feature has
disappeared. The world-line condition \eqref{eq:wlc} was a statement about how
a boost acts on the three spatial components $q_i^b$ of particle $i$ \emph{at
fixed external time $t$}; its inhomogeneous term $t\,\delta^{ab}$ arose
precisely because $t=x_i^0$ was held fixed as an external label while the boost
acted on the remaining three dynamical components. In the SHP formalism, by
contrast, a boost acts homogeneously on \emph{all four} components of $x_i^\mu$
at fixed $\tau$, and $\tau$ itself, being a scalar, is completely unaffected by
the boost. There is therefore no analogue of the inhomogeneous term that made
\eqref{eq:wlc} so rigid, and consequently no constraint forcing the
interaction-dependent piece of the boost generator to track the
interaction-dependent piece of $H$ in the delicate way that
Section~\ref{sec:theorem} showed leads to $\dot{{\bm p}}_i=0$.

Put differently: the CJS/Leutwyler theorem is, at bottom, a statement about
what happens when one tries to single out \emph{one} of the four spacetime
coordinates of an interacting particle and treat it as an external,
non-dynamical parameter while requiring the other three to transform
covariantly. SHP mechanics never makes that choice. All four coordinates of
every particle remain on an equal, fully dynamical footing, and the role that
Newtonian absolute time (or, in the instant form, $t$) played in making
non-relativistic interaction so straightforward is played instead by $\tau$---a
parameter that stands \emph{outside} spacetime altogether rather than being
identified with one of its components. The interaction potential $V$ in
\eqref{eq:shpK} can therefore depend on the full invariant separation
$(x_i-x_j)^2$ and feed directly into a single Hamiltonian for every particle
simultaneously, exactly as an ordinary non-relativistic potential $V(|{\bm
q}_i-{\bm q}_j|)$ depends on the instantaneous relative position under the
Galilei group---because $\tau$ now plays the role that absolute Newtonian time
played there.

Since $p_i^2(\tau)$ is generally $\tau$-dependent, some care is needed to
recover the usual on-shell relativistic particle as a limiting or asymptotic
description. For free particles, or in the asymptotic past/future of a
scattering process where $V\to0$, $p_i^2$ becomes a constant of the
$\tau$-evolution and can be identified with the physical mass-squared $m_i^2$
(up to a controllable spread). More generally, the off-shell mass distribution
can be made sharply peaked around $m_i^2$, so that ordinary relativistic
kinematics is recovered as an excellent approximation whenever the interaction
is weak or the particles are well separated, while genuine off-shell effects
(mass fluctuations induced by the interaction, or by the finite ``width'' with
which $\tau$-evolution explores the mass shell) remain, in principle,
calculable and potentially observable at the level the model is designed to
describe.

The gauge-field companion to this direct many-body potential is the
Stueckelberg--Schwinger--Horwitz--Piron generalization of electrodynamics, in
which a $\tau$-dependent five-dimensional gauge potential $a^\mu(x,\tau)$
couples to a $\tau$-dependent conserved current and mediates the interaction
locally rather than through the direct potential $V$ of \eqref{eq:shpK}; the
resulting off-shell field theory admits both Abelian and non-Abelian versions
and reduces to ordinary Maxwell (or Yang--Mills) theory upon integration over
$\tau$ under suitable conditions \cite{Horwitz2015}. This gauge-field extension
is what makes SHP mechanics a natural setting not only for point-particle
interactions but also for models in which an ``event'' is treated as an
extended object in spacetime, stabilized dynamically rather than kinematically
imposed.


 The overall picture is clear: there is no single ``correct'' way to obtain
 interacting relativistic particle mechanics, only a menu of trade-offs, each
 of which purchases interaction at the cost of one of the three properties that
 make the naive instant-form ansatz so attractive in the first place. Quantum
 field theory purchases it by giving up a closed, finite-dimensional particle
 description; constraint mechanics purchases it by giving up a single universal
 time; covariant action-at-a-distance purchases it by giving up canonical
 structure outright; the practical (Bakamjian--Thomas) Hamiltonian dynamics
 used throughout relativistic few-body physics purchases it by giving up the
 literal identification of the canonical position with a point on a covariant
 world line; and SHP mechanics purchases it by giving up the identification of
 the evolution parameter with one of the particle's own spacetime coordinates,
 restoring instead a genuine fifth, invariant time in the Newtonian sense.

Two general lessons are worth drawing out. First, the no-interaction theorem is
best understood not as an obstacle to be overcome by cleverness within a fixed
framework, but as a precise map of that framework's hidden assumptions: once
the three hypotheses (I)--(III) are identified explicitly, the existence of at
least five structurally distinct resolutions is unsurprising, since each
hypothesis can in principle be relaxed independently, and several ways of
relaxing any one of them may exist. Second, the theorem is a valuable reminder
that ``manifest covariance''---in the strong sense of every dynamical variable
transforming under a fixed, interaction-independent rule---and ``interaction''
are, in a canonical Hamiltonian setting with a single time, close to mutually
exclusive; every viable theory of interacting relativistic particles sacrifices
manifest covariance of \emph{something} (individual positions, a global
Hamiltonian, or the very notion of a fixed number of particle degrees of
freedom) in order to retain the overall, physically observable Poincar\'e
invariance of the theory as a whole.

The no-interaction theorem shows that a specific, seemingly natural way of
writing down relativistic $N$-body Hamiltonian mechanics---canonical variables,
an exact Poincar\'e algebra, and positions that are literal points on covariant
world lines---admits only free motion. This is neither a paradox nor a dead
end: it is a sharply stated structure theorem whose hypotheses can be examined
one at a time and relaxed individually, and the physics literature of the last
sixty years has produced at least five distinct, well-developed ways of doing
so, ranging from the field-theoretic mainstream to the
reparametrization-invariant, off-shell dynamics of Stueckelberg, Horwitz, and
Piron. Far from closing off relativistic particle mechanics, the theorem
supplies the map by which every subsequent construction can be located and
understood.
